% Generated draft: requires mathematical and source review.
\documentclass[11pt]{article}
\usepackage[margin=25mm]{geometry}
\usepackage{fontspec}
\setmainfont{Latin Modern Roman}
\usepackage{amsmath,amssymb,mathtools,unicode-math}
\setmathfont{Latin Modern Math}
\usepackage{enumitem,needspace,xurl,hyperref}
\setlength{\parindent}{0pt}
\setlength{\parskip}{5pt}
\setlength{\emergencystretch}{3em}
\newcommand{\N}{\mathbb{N}}
\newcommand{\Z}{\mathbb{Z}}
\newcommand{\Q}{\mathbb{Q}}
\newcommand{\R}{\mathbb{R}}
\newcommand{\C}{\mathbb{C}}
\newcommand{\F}{\mathbb{F}}
\newcommand{\A}{\mathbb{A}}
\newcommand{\PP}{\mathbb{P}}
\newcommand{\E}{\mathbb{E}}
\newcommand{\eps}{\varepsilon}
\newcommand{\dd}{\,\mathrm{d}}
\begin{document}
\section*{UnsolvedMath ID: 800184}
\textbf{Source problem code:} catalogue entry
\subsection{Definitions and mathematical statement}
Using the conventional interpretation of the notation, let \(d(n)=\#\{a\in\Z_{\geq 1}:a\mid n\}\) be the number of positive divisors of a positive integer \(n\). For real \(x\geq 1\), define
\[
S(x)=\sum_{\substack{n\in\Z_{\geq 1}\\ n\leq x}}d(n)d(n+1)d(n+2).
\]
The problem is to determine an explicit leading asymptotic expression \(A(x)\), positive for all sufficiently large \(x\), such that
\[
S(x)\sim A(x)\qquad(x\to\infty),
\]
where this means \(\lim_{x\to\infty}S(x)/A(x)=1\). The supplied statement does not explicitly define \(d(n)\) or the lower summation boundary; the positive-divisor interpretation and positive-integer summation domain above are conventional interpretations requiring confirmation. No proposed main term, leading constant, or required error estimate is supplied. The source attributes the problem to folklore.
\subsection{Short English statement}
Determine the leading asymptotic behavior, as \(x\to\infty\), of the sum of the products of the divisor counts of three consecutive positive integers. The divisor notation is interpreted conventionally, and the source gives no proposed main term or error bound.
\subsection{Sources}
\begin{itemize}
\item Local UnsolvedMath registry, permanent ID \(800184\), source problem code \texttt{catalogue entry}, title: \emph{Moments of L-functions --- Moments of the Riemann zeta function --- Problem 36}. The supplied registry labels the problem open; this label is not an independently verified mathematical or literature assessment.
\item Hamieh, Kadiri, Ng, \emph{Moments of L-functions Problem List} (2024), Problem 36, \url{https://arxiv.org/abs/2405.17800}. The supplied record lists Hamieh, Kadiri, and Ng under proposer attribution, while the extracted statement itself attributes the problem to folklore.
\item Supplied provenance: accessed 2026-10-16; extraction method \texttt{arxiv\_latex}; upstream revision \texttt{d371ee917ac9647fbb06cd6fc529deec6f3dda82}. These are supplied provenance details; the reference and extraction have not been independently checked here.
\end{itemize}
\end{document}
